\chapter{Methodology}
\label{ch:method}

This chapter describes what was built and how it was evaluated: the dataset and its
verification gate (Section~\ref{sec:data}), how charts are rendered and what that
rendering is allowed to show (Section~\ref{sec:render}), the models
(Section~\ref{sec:models}) and the evaluation protocol (Section~\ref{sec:protocol}).
It closes with the risks that could invalidate the work and how each was handled.

\section{Overview of the Pipeline}

\begin{figure}[h]\centering
\begin{tikzpicture}[
  node distance=0.55cm and 0.75cm,
  box/.style={draw,rounded corners=2pt,align=center,minimum height=1.05cm,
              text width=2.25cm,font=\small},
  ar/.style={-{Latex[length=2mm]},thick}]
\node[box] (a) {Public tables\\+ synthesis};
\node[box,right=of a] (b) {Verification\\gate (tests)};
\node[box,right=of b] (c) {Balanced,\\source-disjoint\\splits};
\node[box,right=of c] (d) {Clean\\rendering};
\node[box,below=1.0cm of d] (e) {Models\\(CNN, VLM)};
\node[box,left=of e] (f) {Evaluation on\\real charts,\\3 renderers};
\node[box,left=of f] (g) {Interactive\\tester};
\draw[ar] (a)--(b); \draw[ar] (b)--(c); \draw[ar] (c)--(d);
\draw[ar] (d)--(e); \draw[ar] (e)--(f); \draw[ar] (e.south) |- ++(0,-0.5) -| (g.south);
\end{tikzpicture}
\caption{Pipeline: tables are admitted only through the verification gate, split by
source, rendered without label-dependent marks, used to train models, and evaluated
on real charts drawn by three plotting libraries.}
\label{fig:pipeline}
\end{figure}

\section{Data: InsightVLM-20K}
\label{sec:data}

\subsection{Task Definition}

\textbf{Input:} a chart image. \textbf{Output:} the statistical finding the chart
shows (one of 20 classes). Every sample additionally carries the data table, the
statistics that support the finding, a reasoning chain and an evidence-grounded
insight, so the same corpus supports table reconstruction, generation and grounding.

\subsection{Label Space}

Twenty finding classes span four analysis families. Each class was admitted only
if it is (i) verifiable by a standard named statistical test, (ii) mutually
exclusive from every other class by explicit guard conditions, and (iii) visible
in a rendered chart.

\begin{table}[h]\centering\small
\caption{Finding families and their classes.}
\label{tab:families}
\begin{tabular}{llp{7.1cm}}
\toprule
\textbf{Family} & \textbf{$n$} & \textbf{Classes} \\
\midrule
Distribution & 6 & normal, right-skew, left-skew, bimodal, heavy-tail, uniform \\
Relationship & 6 & strong positive, strong negative, none, non-linear monotonic,
                  non-linear non-monotonic, heteroscedastic \\
Temporal & 5 & increasing, decreasing, seasonal, changepoint, anomaly spike \\
Group & 3 & significant difference, no difference, dominant category \\
\bottomrule
\end{tabular}
\end{table}

\subsection{The Verification Gate}
\label{sec:gate}

This is the methodological core of the thesis and the mechanism by which
Gap~G1 (Section~\ref{sec:ourgap}) is closed. \textbf{A label is never asserted;
it is tested.} Three rules govern construction.

\paragraph{Rule 1 --- subsample before verifying.} Data is reduced to its final
plotted size \emph{before} the statistical test runs. Verifying a parent
dataframe and then plotting a subsample would produce ground truth describing
data the image does not show.

\paragraph{Rule 2 --- exactly one class must match.} A candidate satisfying two
classes' guards is discarded as ambiguous rather than multi-labelled. This keeps
the primary label single-valued and the confusion matrix interpretable, at the
cost of excluding compound findings (see Section~\ref{sec:risks}).

\paragraph{Rule 3 --- rejection is free; wrong labels are not.} Every guard drops
a sample rather than mislabelling it. More candidates can always be mined; a
poisoned label cannot be undone.

\paragraph{Why robust statistics are mandatory.} A naive implementation of these
tests fails in ways we measured, and three guards exist specifically to prevent
it:

\begin{itemize}[leftmargin=1.3em]
  \item \emph{Moment skewness is not tail-robust.} Over 40 draws of a symmetric
  Student-$t(4)$ sample, the moment estimator alone mislabelled $12.5\%$ as
  skewed. Requiring agreement with Bowley (quartile) skewness eliminated this.
  \item \emph{Hartigan's dip test is calibrated against the uniform}, so a
  genuinely flat distribution trips $p<0.05$ roughly $5\%$ of the time. The
  bimodal guard therefore additionally requires a KS test against the uniform to
  reject.
  \item \emph{A monotone ramp defeats changepoint detection}, producing a
  high-SNR ``changepoint''. The guard therefore requires a two-level step model to
  beat a linear fit \emph{and} an abruptness ratio $>0.55$ (measured: true step
  $1.15$, smooth rise-then-fall $0.18$, pure ramp $0.14$).
\end{itemize}

Generator-to-class fidelity over 10 seeds $\times$ 20 classes is $95.5\%$; every
remaining miss is a conservative rejection, never a wrong label.

\subsection{Two Construction Tracks}

A balanced corpus cannot be built from real data alone. Public repositories hold
thousands of right-skewed columns and almost no clean inverted-U relationships:
mining 719 real tables yielded exactly \textbf{ten}. The dataset is therefore a
declared hybrid, and we state the mix rather than obscuring it.

\emph{Track R (real-seeded)} mines every plottable view from 719 public tables
(OpenML, World Bank, Our World in Data, vega-datasets, seaborn-data, UCI).
\emph{Track S (controlled synthesis)} generates data whose property holds by
construction, from a domain vocabulary spanning health, economics, environment,
engineering, education and retail. \textbf{Both tracks pass the identical gate.}

Provenance is balanced \emph{within} each class rather than real-first;
real-first would leave some classes almost entirely real and others almost
entirely synthetic, letting a model key on provenance instead of the statistical
pattern.

\subsection{Splits and the Supplementary Real Pool}
\label{sec:splits}

Splits are $14{,}000/2{,}000/4{,}000$, exactly $700/100/200$ per class, and
\textbf{source-disjoint} for real data: a hash of the source-table identifier
assigns every chart derived from one table to a single split. The real subsets hold
$4{,}338$ training, $553$ validation and $1{,}192$ test charts; real class frequencies
are uneven because public data are.

Balance caps real data at 500 per class, so the gate admitted more real charts than
the balanced splits use. We keep the surplus as a \textbf{supplementary real training
pool}: candidates whose source table the same hash assigns to \emph{train}, with exact
duplicates of any dataset sample removed and at most 400 per class. The build asserts
that no validation or test source contributes. The pool adds $2{,}545$ real training
charts in nine classes, raising real training data from $4{,}338$ to $6{,}883$.

\begin{table}[h]\centering\small
\caption{Dataset summary and integrity audit.}
\label{tab:dataset}
\begin{tabular}{lr}
\toprule
Samples & $20{,}000$ \\
Classes & 20, exactly $1{,}000$ each \\
Chart types & 26 \\
Provenance & $6{,}083$ real (654 source tables) / $13{,}917$ synthetic \\
Supplementary real training pool & $2{,}545$ (train-split sources only) \\
Underlying data points & $8{,}487{,}134$ \\
\midrule
Label re-verification from stored data & $100.00\%$ \\
Leaked sources across splits & 0 of 654 \\
Duplicate data fingerprints & 0 \\
Real test charts with near-duplicate data in training (Sec.~\ref{sec:rq5}) & 85 of 1{,}192 \\
Grounding boxes outside canvas & 0 \\
\bottomrule
\end{tabular}
\end{table}

\section{Rendering}
\label{sec:render}

Charts are drawn with matplotlib and seaborn on a fixed $768\times576$ canvas across
26 chart types and 8 visual styles. Each sample's style is chosen from a CRC32 hash of
its identifier, so rendering is reproducible across processes.\footnote{The first
version used Python's built-in \texttt{hash()}, which is salted per process, so styles
could not be reproduced. This was corrected in version 2.} Grounding boxes marking the
evidence (outlier band, modal bin, change point, anomalous points, dominant wedge) are
emitted in image pixel coordinates. Tight bounding-box cropping is deliberately
disabled: it resizes the canvas \emph{after} the axes transform is read, silently
invalidating every box. Dates stored as epoch timestamps are plotted as calendar dates.

\subsection{No Label-Dependent Marks (Version 2)}
\label{sec:v2}

The first release drew annotations on time-series charts that depended on the label:
a fitted trend line only on increasing and decreasing series, a vertical line at the
change point only on level-shift series, and circles only around anomalies. The
seasonal class carried no mark, so the absence of a mark was itself informative. Such
marks let a model identify the class without reading the data. Version 2 draws no
label-dependent marks; the evidence survives only as grounding boxes in the metadata.
No other family drew label-dependent marks (regression lines on scatter plots are drawn
independently of the label). Chapter~\ref{ch:results} measures how strongly models had
exploited the marks. All results in this thesis use version 2.

\subsection{Cross-Renderer Test Sets}

All training charts come from one plotting library, so a model could learn
matplotlib's appearance rather than the statistics. The 1{,}192 real test tables were
therefore re-drawn by two more libraries from the same stored data:
\emph{Plotly}, with five templates (excluding the one that imitates seaborn), and
\emph{Vega-Lite}, with nine themes (excluding those that match a matplotlib training
style). Chart types map to each library's nearest native form; where none exists the
substitute is a form the taxonomy lists as valid for every class that chart type can
carry, so no label becomes unreadable.

\subsection{Training-Only Renderer Augmentation}

Every training and supplementary chart was also rendered twice more, for training
only: by Plotly, and by matplotlib with randomised fonts, font sizes, line and marker
sizes, colours, background tints, gridlines, spines and title placement. The data and
chart type are unchanged, so the label is too. No validation or test chart is rendered
this way and no Vega-Lite render is used in training, so \textbf{Vega-Lite remains a
renderer that no model has seen}, and the Plotly test set becomes a seen-renderer test
for augmented models.

\section{Evaluation Protocol}
\label{sec:protocol}

This section addresses RQ5 and objective O6.

\subsection{Why the Conventional Baseline Is Wrong Here}

With 20 balanced classes, chance is $5\%$. But chart type is associated with the
finding \emph{by construction} --- a pie chart can only carry a composition
finding, a QQ plot only a distribution finding. A model consulting \emph{no pixels
at all}, predicting the most common class for each chart type, reaches
$\mathbf{0.288}$. Reporting against $5\%$ would credit a model with roughly 24
points it did not earn. We therefore use the chart-type prior as the floor.

\subsection{A Data Oracle Gives the Ceiling}

Running the verification battery on each sample's stored data table --- seeing the
numbers, not the picture --- yields $\mathbf{0.9995}$. This is the deterministic
rule-based pipeline of RQ2/O3, and a vision model is interesting only within
$[0.288,\ 0.9995]$.

\subsection{Real Charts Are the Headline}

Because the corpus is majority synthetic, mixed-test accuracy cannot distinguish
learning the statistics from learning the generator. The headline metric is accuracy
on the $1{,}192$ real test charts (RQ4/O5), reported for the matplotlib render and for
the Plotly and Vega-Lite re-renders. Every accuracy carries a 95\% bootstrap confidence
interval (4{,}000 resamples), and differences between systems scored on the same
charts are tested with exact McNemar tests. \textbf{Every choice of model, checkpoint,
ensemble member and hyperparameter is made on the validation split}; the test split is
scored once.

\subsection{Shortcut and Bias Audit}

Eleven failure modes are measured against thresholds fixed in advance
(Table~\ref{tab:bias}). All pass on version 2.

\begin{table}[h]\centering\small
\caption{Shortcut and bias audit on version 2; all eleven tests pass.}
\label{tab:bias}
\begin{tabular}{lrr}
\toprule
\textbf{Test} & \textbf{Value} & \textbf{Threshold} \\
\midrule
Class balance (imbalance ratio) & 1.000 & $=1.000$ \\
Split balance (max s.d.) & 0.0000 & $<0.01$ \\
Chart type $\rightarrow$ label (pixel-free prior) & 0.2883 & reported as floor \\
Render style $\rightarrow$ label (Cram\'er's $V$) & 0.0000 & $<0.05$ \\
\textbf{Provenance detectability (AUC)} & \textbf{0.551} & $<0.75$ \\
Axis-text shortcut (macro-F1) & 0.3688 & $<0.55$ \\
Pixel-statistic shortcut (macro-F1) & 0.4038 & $<0.55$ \\
Source concentration (top share) & 17.2\% & $<35\%$ \\
Length confound (within family) & 2.04 & $<4.0$ \\
Subject-domain balance & 22.8\% & $<45\%$ \\
Category-vocabulary spread & 13.0\% & $<35\%$ \\
\bottomrule
\end{tabular}
\end{table}

Provenance detectability matters most: with $69.6\%$ of samples synthetic, the obvious
objection is that a model learns to recognise generated charts. A classifier on
low-level pixel statistics separates real from synthetic at $\mathrm{AUC}=0.551$ ---
close to chance. The audit examines metadata and pixel statistics, however, not what a
trained model attends to; it did \emph{not} detect the version-1 marks
(Section~\ref{sec:v2}), which were found by inspecting the renderer. We therefore added a
direct ablation (Section~\ref{sec:leak}).

\section{Models}
\label{sec:models}

\subsection{Discriminative Models (O2)}

ImageNet-pretrained ResNet-101, ResNet-50 and DenseNet-121 are fine-tuned at
$384\times384$ for up to 12 epochs with AdamW (learning rate $3\times10^{-4}$,
weight decay 0.05), a one-cycle schedule, label smoothing 0.1 and mixed precision on
two T4 GPUs. Two choices target real-chart accuracy directly:
\begin{itemize}[leftmargin=1.3em]
  \item \textbf{Real-data oversampling.} A weighted sampler draws about half of every
  batch from real charts (31\% of the balanced training split is real).
  \item \textbf{Real-validation checkpointing.} The checkpoint kept is the epoch with
  the highest accuracy on the 553 \emph{real} validation charts, with mixed validation
  accuracy as the tie-break. Mixed accuracy saturates near 0.95 and says little.
\end{itemize}
Models are trained in three steps, each an ablation of the next: on the balanced split;
with the supplementary real pool added; and with renderer augmentation added (each
epoch shows each training chart as one of its three renders, uniformly).

\paragraph{Augmentation constraint.} Horizontal and vertical flips are
\emph{excluded}. Mirroring a scatter converts a strong positive correlation into
a strong negative one, and mirroring a line chart converts an increasing trend
into a decreasing one --- both separate classes here, so the default augmentation
would silently corrupt roughly $30\%$ of labels. Only label-preserving transforms
(mild crop and scale, colour jitter, blur, small erasures) are used.

\paragraph{Ensemble.} The deployed model averages the class probabilities of the
full-recipe members (supplementary pool and renderer augmentation) with equal weights.
Probabilities are averaged rather than logits, whose scales differ across
architectures. A label-shift correction towards real-data class frequencies, estimated
from training data only, is available with its strength tuned on validation (it is
not used: the tuned strength is zero). Greedy ensemble selection on validation and the
average of all trained models are reported as alternatives.

\subsection{Zero-Shot Vision--Language Models}

To test whether current general-purpose VLMs can already do the task, Qwen2.5-VL-3B,
Qwen2.5-VL-7B and Qwen3-VL-8B receive each real test chart with all 20 class names and
one-line definitions and must answer with one class name (greedy decoding, no task
training). Answers are parsed by exact label match; over 99.9\% parse.

\subsection{Generative Model: InsightVLM (O2, O3)}

Qwen2-VL-2B adapted with QLoRA ($r=16$, 4-bit NF4, $18.5$M trainable parameters,
$0.83\%$ of the model), trained one epoch on $4{,}000$ samples. The supervision
target is the reasoning chain plus the insight, so the model learns to produce
evidence-grounded prose rather than a class identifier. Numeric fidelity is the share
of numbers the model cites that match the ground-truth statistics to two decimals.
This model was trained and evaluated on version-1 renders before the leak was found;
Chapter~\ref{ch:results} reports it with that caveat.

\subsection{Evidence-Grounding Evaluation (O4)}

For each real test chart that carries at least one evidence box, the Grad-CAM map of
each deployed member for its predicted class is computed at the last convolutional
block, upsampled, and its peak mapped back to image pixels. A \emph{hit} is a peak
inside any evidence box. The control is a uniformly random point inside the plot area,
whose expected hit rate equals the evidence boxes' share of the plot area, so large
evidence regions cannot inflate the result.

\subsection{Interactive Tester}

A Gradio web application (Appendix~\ref{app:tester}) loads the deployed ensemble and
runs uploaded charts through exactly the pipeline that was scored. It reports the
finding, its confidence, how many members agree, and the statistical test that defines
the class, and can draw random held-out test charts to show predictions against the
verified ground truth.

\section{What Could Go Wrong}
\label{sec:risks}

\begin{description}[leftmargin=1.4cm,style=nextline,font=\bfseries]
  \item[R1. Synthetic majority.] $69.6\%$ of samples are synthetic; balance forced
  this. \emph{Mitigation:} provenance balanced within each class, identical gate
  for both tracks, real charts reported as the headline, and a supplementary real pool
  for training. \emph{Residual risk:} the real subset may not represent real charts
  broadly.

  \item[R2. Single-label design.] Compound findings were excluded, not handled, so
  the model will not learn to report a trend and an anomaly together.
  \emph{Mitigation:} stated as a limitation; multi-label extension is future work.

  \item[R3. Difficulty axis does not track difficulty.] Accuracy does not follow
  the evidence-margin tertiles (Spearman $0.034$). Extreme margins produce
  degenerate charts --- a residual spread ratio of $1537$ renders as a single
  vertical smear. \emph{Mitigation:} reported as a finding about chart benchmarks
  rather than presented as a working difficulty axis.

  \item[R4. Renderer monoculture.] Training charts come from one plotting library.
  \emph{Mitigation:} Plotly and Vega-Lite test sets measure the effect, and
  training-only renderer augmentation reduces it; Vega-Lite is held out entirely.

  \item[R5. Thresholds are conventions.] $|r|>0.7$ as ``strong'' and
  $\varepsilon^2>0.14$ as ``large'' are choices. \emph{Mitigation:} all cuts are
  recorded explicitly so results can be re-derived under alternatives.

  \item[R6. Template recitation by the VLM.] The training target embeds each class
  description verbatim, so the model could score well by memorising a phrase.
  \emph{Mitigation:} the real/synthetic gap is the diagnostic --- observed
  $0.683$ vs $0.844$, comparable in magnitude to the CNNs, indicating chart
  reading rather than recitation.

  \item[R7. Compute constraints.] Training ran on free Kaggle T4 GPUs with a weekly
  quota; the generative result used $4{,}000$ samples, one epoch and 2B parameters at
  4-bit. \emph{Mitigation:} reported as a compute-constrained floor, not a ceiling.

  \item[R8. Rendering leaks.] A renderer can reveal the label without the audit
  noticing. This happened (Section~\ref{sec:v2}). \emph{Mitigation:} version 2 removes
  every label-dependent mark, and a with/without ablation quantifies the effect.

  \item[R9. Selection on the test set.] An earlier ensemble was chosen by ranking
  combinations on test accuracy, which inflates the reported number.
  \emph{Mitigation:} all selection now uses validation only (Section~\ref{sec:protocol}).
\end{description}

\section{Reproducibility}

The corpus rebuilds end to end with \texttt{run\_all.py}; \texttt{07\_validate.py}
(integrity) and \texttt{10\_bias\_audit.py} (shortcuts) audit it. The clean renders,
the supplementary pool, the Plotly and Vega-Lite renders and the augmentation renders
are produced by \texttt{06\_render.py}, \texttt{12\_extra\_real.py},
\texttt{11\_render\_plotly.py} and \texttt{13\_render\_vega.py}. Training notebooks are
generated by \texttt{train/make\_split\_notebooks.py}, and every reported number comes
from \texttt{final\_ensemble.py}, \texttt{stats\_report.py},
\texttt{grounding\_eval.py} and \texttt{ablate\_marks.py}. All thresholds live in a
single module so results can be re-derived under alternative conventions.
Documentation follows the datasheet standard of \citet{datasheets}.
